\section{Geometric components over a field of Weber values}

% BEGIN W76F1
\paragraph{W76F1. Full torsion on a conjugate CM elliptic curve.}
Write $s=H.\mathrm{rt0}$ and $A=\mathbb Z[(1+s)/2]$ in the
chosen complex embedding. It is the integer ring of
$k=\mathbb Q(\sqrt{-d})$, with $d$ the host prime,
$d\equiv7\pmod8$. Let $F\subset\mathbb C$ be the given coefficient
field containing $s$. For positive $N$, the native condition
$\mathrm{WeberIn}(H,F,N)$ implies $\mathrm{CurveOver}(H,F,N)$.
This last assertion allows the comparison isomorphism of the
elliptic curve to be chosen for each field automorphism fixing $F$.

In fact take the elliptic curve $E_0=\mathbb C/A$ in
Weierstrass coordinates, write $\Delta=g_2^3-27g_3^2$, and
set $h(z)=(g_2g_3/\Delta)\wp(z)$. The supplied $g_2^3/\Delta$
determines its $j$-invariant. Neither $g_2$ nor $g_3$ vanishes:
otherwise the Weierstrass equation has an elliptic automorphism
of order three or four fixing the origin. But the endomorphisms
of $\mathbb C/A$ fixing the origin are multiplication by the
multipliers of its lattice, namely $A$, and $A^\times=\{\pm1\}$.
For $\tau\in\operatorname{Aut}(\mathbb C/F)$ equality of the
$j$-invariants gives an origin-preserving isomorphism between
the conjugate and original curves. The Weber function is carried
to its original normalization: scaling coordinates by
$(u^2,u^3)$ scales $(g_2,g_3,\Delta)$ by $(u^4,u^6,u^{12})$,
so $h$ is unchanged. Its fibers away from zero on the curve
consist of $\{P,-P\}$ because $\wp$ is the $x$-coordinate.

The Weber hypothesis fixes $h(P)$ for every nonzero
$P\in E_0[N]$. Thus under this comparison each $P$ goes to $P$
or $-P$. The sign can be made simultaneous. Choose a point sent
to $-P\ne P$, if one exists. If $Q$ is fixed, the two signs
allowed for $P+Q$ give $2P=0$ or $2Q=0$ on comparing
$-P+Q$ with $\pm(P+Q)$. The first case is excluded, so $Q=-Q$.
Composition with inversion now fixes the entire $N$-torsion.
The zero point and $N=1,2$ satisfy the same conclusion. The
comparison intertwines the $A$-action. The action on differentials
of multiplication by $a\in A$ is the scalar $a$ in the chosen
embedding, fixed by $\tau$. In characteristic zero this scalar
determines the endomorphism (lift to the complex universal cover).
We have obtained exactly the group pullback with trivial torsion
multiplier in $\mathrm{FixesCurve}$. If $N'\mid N$, every
$N'$-division point is an $N$-division point, so the same
comparison gives $\mathrm{CurveOver}(H,F,N')$.
% END W76F1

% BEGIN W76F2
\paragraph{W76F2. A split period in each complex component at every allowed level.}
In an original pinned lower row let $p,q$ be the signature,
$r=p+q$, and let $E/E^+$ be the host CM extension. The same
argument applies to every component lattice at every level
$l'$ on this row. There exists an original polarized period on
that lattice splitting it as a direct sum of a pure $s$ part
and a pure $-s$ part over $A$. In particular the corresponding
abelian variety is a direct factor of a power of the elliptic
curve of W76F1.

Here is why an \emph{integral} decomposition is available.
The rank-one host datum of signature $(1,0)$ is carried by the
fractional ideal $\mathfrak A$ of PCL8: take its alternating trace
form with the sign of $-\xi_H$. The product with $\bar{\mathfrak A}$
and the different gives the integral dual equal to $\mathfrak A$;
the signed imaginary embedding gives the indicated pure type.
Negating the form supplies the opposite datum. Their orthogonal
sum in $p$ and $q$ copies gives a split self-dual reference lattice
$L_0$ with rational summands $V_+,V_-$. PUN3--PUN5 put a given
component lattice $L$ of this signature in the same rational
Hermitian space and give a local integral isometry $L_{0,v}\to L_v$
at every finite place of $E^+$. In particular the conclusion at
dyadic places uses the host's unramified integral duality there.

If either block is empty the original lattice is already pure.
Otherwise let $P$ be the orthogonal projector onto $V_+$.
For every finite $v$ choose such a local isometry $g_v$,
taking the identity wherever $L_0=L$ locally. The conjugate
$g_v P g_v^{-1}$ is integral on $L_v$. Correct the determinant
of $g_v$ by composing on the right with a unitary automorphism
commuting with $P$: multiply a nondegenerate line in $V_+$ by
$\det(g_v)^{-1}$ and fix its orthogonal complement. The determinant
has norm one; a rational line of nonzero Hermitian norm in
$V_+$ works at every place, including the split ones. This
correction need only preserve the projector. Thus the resulting
local operators lie in $\mathrm{SU}(V)$, with the same integral
conjugate projector. The condition $uPu^{-1}L_v\subset L_v$ is
open. At almost all $v$ it contains the integral stabilizer.
The special unitary group over $E^+$ is simply connected and
absolutely almost simple (here the pinned rank is at least six).
At an archimedean place with both signs it is noncompact.
By the strong approximation theorem its rational points are dense
in its finite adelic group (Platonov--Rapinchuk, \emph{Algebraic
Groups and Number Theory}, Theorem 7.12). Hence some single
rational $u$ has $uPu^{-1}$ integral on $L$ at every finite
place. Intersection of the local lattices gives integrality
on $L$. The image and kernel supply orthogonal complementary
integral modules, with their original definite signs at the
real places.

On their real spans take the pure complex structures belonging
to the two conjugate embeddings $k\hookrightarrow\mathbb C$.
The signs just proved are the Riemann positivity condition and
orthogonality makes the cross term vanish. Choosing bases of the
complex tangent eigenspaces produces a period in the original
signature domain; one can order its coordinates by the prescribed
CM multiplicities. This choice preserves the entire given lattice
and its level marking. Each signed lattice regarded over the
Dedekind ring $A$ is finite projective. With the action on the
negative one conjugated, write their sum as a direct summand
of $A^m$. The idempotent in $M_m(A)$ gives a direct factor
of $E_0^m$ with precisely the just constructed period lattice
and complex structure. This uses an integral summand, so no
isogeny denominator is introduced into the level.
% END W76F2

% BEGIN W76F3
\paragraph{W76F3. Fixed components and their constants at the prescribed level.}
Take any upper level $l'$ on an original pinned row, any
corresponding moduli scheme $X'_o$ and coefficient record $c'$.
Put $F=c'.F$. If $N>0$, $\operatorname{lcm}(l'.M,l'.A)\mid N$
and $\mathrm{WeberIn}(H,F,N)$ or $\mathrm{CurveOver}(H,F,N)$
holds, then $\mathrm{FnAlgClosed}(X'_o,c')$ holds. Explicitly,
for each irreducible component $Z$ of the model over $F$,
every element of the original framed function field $\Omega(Z)$
algebraic over $F$ belongs to $F$.

The host root belongs to the fixed base subfield of any such
$c'$. In a supplied curve comparison the marked elliptic lattice
is a complex homothety of $A$; the homothety identifies the
curve with $E_0$ together with all its labeled division points.
Use W76F1 or this supplied comparison, and restrict
to the level $N'=\operatorname{lcm}(l'.M,l'.A)$. For each
automorphism of $\mathbb C/F$, its comparison on $E_0$ carries
to every power and all its $A$-linear direct factors. More
explicitly, the $A$-linear projectors, inclusions and projections
are matrices of elements of $A$ and hence intertwine the
comparison. Via the canonical principal polarization of $E_0$
the same holds for homomorphisms to the dual. Every homomorphism
between the factors or their duals extends by projection and
inclusion to such a matrix. Thus the induced comparison preserves
the $\mathcal O_E$-action and the particular polarized homomorphism
of W76F2. On an integral summand the $N'$-torsion is the image
under its idempotent of $E_0^m[N']$. The comparison fixes its
labeled torsion, in particular both original level markings.
Equivalently in period coordinates one obtains an integral
isometry with the two original marking congruences; the
marked-period criterion PUN23 identifies the moduli points.
The normalized symmetric line in that criterion is determined by
the prescribed host semicharacter of the alternating form. This
argument uses it on the same integral lattice.

Every geometric component contains such a point: the actual
marked period family covers the complex model, and a connected
period domain containing a preimage can be moved within its
domain to the split period of W76F2 on its own lattice. The
domain is connected by NFR2, and smoothness makes irreducible
components disjoint open and closed. A field automorphism fixing
$F$ permutes these algebraic components; it fixes the split
moduli point just constructed and therefore fixes its component.
This reasoning applies separately to every component.

For completeness, a smooth finite-type variety over $F$ with this
property has geometrically integral irreducible components. For
one irreducible $Z$ the absolute Galois group acts transitively
on the components of $Z_{\bar F}$: a proper nonempty orbit union
would descend to a proper open and closed subset of $Z$, because
after smooth base change the components are disjoint. The components
are defined over finite extensions, and each automorphism of
$\bar F/F$ extends to $\mathbb C$ by fixing a transcendence
basis and extending algebraically. Base extension from $\bar F$
to $\mathbb C$ preserves irreducible components. The preceding
fixedness makes each orbit a singleton, so $Z_{\bar F}$ is
integral. A proper finite extension $F'/F$ in $\kappa(Z)$
would give nontrivial idempotents in
$F'\otimes_F\bar F\subset\kappa(Z)\otimes_F\bar F$, contradicting
integrality, since the characteristic is zero. Finally
$\Omega(Z)=\kappa(Z)(t_1,\dots,t_d)$ in the original frame
variables. A field is relatively algebraically closed in a purely
transcendental extension (a nonconstant rational function in one
variable is transcendental, and induction removes the variables).
This proves the claimed assertion on $\Omega(Z)$ for the actual
upper model, including the zero-dimensional case.
% END W76F3
